% Chapter 3 — Canonicalization and invariant front-ends
% G4 content: §3.6 deep phase-invariant encoding (CLAIMS row 15) written; §§3.1-3.5 stubs
% pointing at the implemented canonicalizers, to be filled at G5.
\documentclass[11pt]{article}
\usepackage{amsmath,amssymb,amsthm,mathtools}
\newtheorem{theorem}{Theorem}[section]
\newtheorem{proposition}[theorem]{Proposition}
\newtheorem{lemma}[theorem]{Lemma}
\newtheorem{corollary}[theorem]{Corollary}
\newtheorem{definition}[theorem]{Definition}
\newtheorem{remark}[theorem]{Remark}
\newcommand{\R}{\mathbb{R}}
\newcommand{\Z}{\mathbb{Z}}
\title{Chapter 3: Canonicalization and invariant front-ends}
\begin{document}
\maketitle

\section{Overview}
Chapter~1 established that the per-neuron symmetry group of a sine network is $D_\infty$ and that
the per-layer group is $D_\infty\wr S_n$, with identifiability proved at $L=1$. This chapter asks
what can be \emph{done} with that structure: exact canonicalization ($c_{\mathrm{sort}}$,
$c_{\mathrm{align}}$), its impossibility in continuous form (PO-5), and invariant encodings that
bypass canonicalization altogether. Sections~\ref{sec:csort}--\ref{sec:margins} are written at G5;
Section~\ref{sec:deep} is written here because the $L=2$ construction it contains was needed to
run rung W10 of Study~S1.

\section{The template-free canonicalizer $c_{\mathrm{sort}}$}
\label{sec:csort}
[G5] Phase reduction to $b\in[-\pi/2,\pi/2)$, sign fixing against a reference direction, and
lexicographic ordering of neurons; exactness with respect to the implemented group; tie margins.

\section{The template canonicalizer $c_{\mathrm{align}}$}
\label{sec:calign}
[G5] Activation matching against a template with a Hungarian assignment, the correlation sign as
the $\sigma$ fix, and function preservation asserted on every call. Empirically this is the
strongest rung of the S1 ladder that is not an oracle (Chapter~5), and the template it needs is
the chapter's main practical caveat.

\section{Margins, ties, and the continuity obstruction}
\label{sec:margins}
[G5] PO-5's winding argument; where real fits land relative to tie sets (figure F7).

\section{Deep phase-invariant encoding}
\label{sec:deep}

Proposition PO-4 gives a family of $D_\infty$-invariants of a single neuron $(w,b,u)$, organised by
the parity of each generator's action. For a network of depth $L=1$ that family is complete enough
to separate orbits on the generic stratum. It does not, however, transfer to $L\ge 2$: the outgoing
vector $u$ of a hidden neuron is a \emph{column of the next layer's weight matrix}, and the next
layer's own group acts on it. An encoding built from $\sin b\,u$ at layer~1 is therefore invariant
under layer~1's action and covariant under layer~2's --- not invariant at all. This section closes
the gap for $L=2$.

\subsection{Setup and notation}
Write the stored (canonical) form of a two-layer sine network as
\begin{equation}
h_1 = \sin(W_1x+b_1),\qquad h_2=\sin(W_2h_1+b_2),\qquad f = W_3h_2+b_3,
\end{equation}
with $W_1\in\R^{n\times m}$, $W_2\in\R^{p\times n}$, $W_3\in\R^{c\times p}$. For a layer-1 neuron
$i$ put $w_i = (W_1)_{i:}$, $b_i=(b_1)_i$ and $u_i = (W_2)_{:i}$; for a layer-2 neuron $k$ put
$w_k'=(W_2)_{k:}$, $b_k'=(b_2)_k$ and $u_k'=(W_3)_{:k}$. By Theorem PO-1 the group acts per neuron
by $g_{d,j}:(w,b,u)\mapsto((-1)^dw,\,(-1)^db+\pi j,\,(-1)^{d+j}u)$, together with a permutation of
the neurons of each layer.

Throughout, write $\varepsilon_i := (-1)^{d_i+j_i}$ for the scalar that layer~1's action at
neuron~$i$ applies to $u_i$. The parity bookkeeping we need is elementary:
\begin{equation}
\label{eq:parities}
\begin{aligned}
w_i &\mapsto (-1)^{d_i}w_i, &\qquad \cos 2b_i &\mapsto \cos 2b_i, &\qquad \|w_i\|^2&\mapsto\|w_i\|^2,\\
\sin 2b_i &\mapsto (-1)^{d_i}\sin 2b_i, &\qquad \sin b_i&\mapsto \varepsilon_i \sin b_i,
&\qquad \cos b_i &\mapsto (-1)^{j_i}\cos b_i .
\end{aligned}
\end{equation}
Only the third and fifth entries are not immediate: $\sin((-1)^db+\pi j) = (-1)^j\sin((-1)^db)
= (-1)^{j}(-1)^{d}\sin b = \varepsilon\sin b$, and $\cos((-1)^db+\pi j)=(-1)^j\cos b$.

\subsection{The layer-2 Gram matrix}

\begin{lemma}[Gram invariance]
\label{lem:gram}
Let $G := W_2^{\top}W_2\in\R^{n\times n}$, i.e.\ $G_{i\ell} = \langle u_i,u_\ell\rangle$. Then $G$ is
invariant under the entire layer-2 group, and under layer~1's action it transforms as
$G_{i\ell}\mapsto \varepsilon_i\varepsilon_\ell G_{i\ell}$, i.e.\ $G\mapsto PDGDP^{\top}$ with
$D=\mathrm{diag}(\varepsilon)$ and $P$ the layer-1 permutation.
\end{lemma}

\begin{proof}
A layer-2 element acts on neuron $k$ by $w_k'\mapsto(-1)^{d_k}w_k'$, $b_k'\mapsto(-1)^{d_k}b_k'+\pi
j_k$, $u_k'\mapsto(-1)^{d_k+j_k}u_k'$, and by permuting the layer-2 neurons. In terms of $W_2$ the
first is a sign flip of \emph{row} $k$ and the last is a permutation of rows; the bias and $W_3$ are
untouched by $G$. Each entry of $G$ is a sum over rows,
$G_{i\ell}=\sum_k (W_2)_{ki}(W_2)_{k\ell}$, in which a row sign $(-1)^{d_k}$ occurs twice and
cancels, and a row permutation only reindexes the sum. Hence $G$ is layer-2 invariant. Layer~1's
action multiplies the column $u_i$ by $\varepsilon_i$, giving the stated covariance, and permutes
the columns, giving the conjugation by $P$.
\end{proof}

\subsection{Sign-cancelling couplings}

\begin{proposition}[Invariant matrices]
\label{prop:matrices}
Define, for layer-1 neurons $i,\ell$,
\begin{align}
A_{i\ell} &:= \sin b_i \sin b_\ell\, G_{i\ell}, &
B_{i\ell} &:= \langle \cos b_i\, w_i,\ \cos b_\ell\, w_\ell\rangle\, G_{i\ell}, \notag\\
C_{i\ell} &:= \langle \sin 2b_i\, w_i,\ \sin 2b_\ell\, w_\ell\rangle . &&
\end{align}
Each is symmetric, invariant under the layer-2 group, and satisfies $M\mapsto PMP^{\top}$ under the
layer-1 group. Consequently the sorted spectrum of each is invariant under the full group
$\prod_{l}\left(D_\infty^{n_l}\rtimes S_{n_l}\right)$.
\end{proposition}

\begin{proof}
Layer-2 invariance is Lemma~\ref{lem:gram} for $A,B$ (the layer-1 quantities $w,b$ are untouched by
layer~2) and is vacuous for $C$. For layer~1, use \eqref{eq:parities}. In $A$, the factor
$\sin b_i\sin b_\ell$ acquires $\varepsilon_i\varepsilon_\ell$ and $G_{i\ell}$ acquires
$\varepsilon_i\varepsilon_\ell$, whose product is $1$. In $B$, the vector $\cos b_i w_i$ acquires
$(-1)^{j_i}(-1)^{d_i}=\varepsilon_i$, so the inner product acquires
$\varepsilon_i\varepsilon_\ell$, again cancelled by $G_{i\ell}$. In $C$, each factor
$\sin 2b_i\,w_i$ acquires $(-1)^{d_i}(-1)^{d_i}=1$ and is invariant on its own. Permuting layer-1
neurons permutes indices $i,\ell$ simultaneously, which is conjugation by $P$; the spectrum of a
symmetric matrix is invariant under conjugation by a permutation matrix and under the diagonal sign
conjugation already cancelled, so sorting the eigenvalues yields group invariants.
\end{proof}

\begin{remark}
The pairing is forced, not chosen. A quantity with layer-1 parity $\varepsilon$ cannot be made
invariant by itself, because $\varepsilon_i$ depends on the neuron; it must be contracted against
another $\varepsilon$-covariant quantity carrying the \emph{same} index, and $G$ is the only object
available that is simultaneously layer-2 invariant and $\varepsilon_i\varepsilon_\ell$-covariant.
This is why $\cos b_i\cos b_\ell G_{i\ell}$ --- the apparent companion to $A$ --- is \emph{not}
invariant: $\cos b$ carries $(-1)^{j}$ rather than $\varepsilon$, leaving a residual
$(-1)^{d_i+d_\ell}$.
\end{remark}

\begin{proposition}[Order statistics]
\label{prop:order}
The per-neuron scalars $\|w_i\|^2$, $\cos 2b_i$ and $G_{ii}=\|u_i\|^2$ are invariant under the
per-neuron action and are permuted by the layer-1 permutation. Their order statistics --- and, more
usefully, their values gathered under a single common sort key that is itself invariant --- are
therefore full-group invariants.
\end{proposition}

\begin{proof}
Immediate from \eqref{eq:parities} and Lemma~\ref{lem:gram} (diagonal entries have
$\varepsilon_i^2=1$). Sorting by a common invariant key preserves the association between the three
features, which independent sorting would destroy.
\end{proof}

\subsection{The encoding, and what it is not}

\begin{definition}[Deep invariant encoding]
For an $L=2$ sine network let
\begin{multline}
\Phi_{\mathrm{deep}} := \bigl(\mathrm{spec}(A),\ \mathrm{spec}(B),\ \mathrm{spec}(C),\\
\text{key-sorted } (\|w_i\|^2,\cos 2b_i, G_{ii})_{i\le n},\
\Phi_{\mathrm{PO4}}(\text{layer }2)\bigr),
\end{multline}
where the layer-2 block applies the PO-4 invariants to $(w_k',b_k',u_k')$ contracted over the
layer-1 index (so that only $\|w_k'\|^2$, $\cos 2b_k'$, $\|u_k'\|^2$ and $\sin b_k'\,u_k'$ enter,
each of which is insensitive to layer~1's action on the column index).
\end{definition}

\begin{corollary}
$\Phi_{\mathrm{deep}}$ is invariant under the full group. For the frozen corpus geometry
($n=p=32$, $m=2$, $c=1$) it has dimension $320$.
\end{corollary}

Invariance is certified numerically as well as proved: applying random group elements with windings
up to $3$ and non-trivial permutations moves $\Phi_{\mathrm{deep}}$ by $\approx 3\times 10^{-7}$
relative, which is \texttt{float32} round-off (test T10, which also verifies that the encoding is
not trivially constant across distinct networks).

Two limitations should be stated plainly. First, $\Phi_{\mathrm{deep}}$ is invariant but is
\emph{not claimed to be separating}: spectra discard the eigenvectors, and Chapter~5 finds that
rung W10 reaches $35.5\%$ where the template canonicalizer reaches $64.4\%$, which is consistent
with a real loss of information in the pooling (though it does not prove one, since the two rungs
also differ in dimension). Second, the construction is specific to $L=2$. At $L\ge 3$ one needs a
Gram per successive layer and the parity bookkeeping compounds; whether a finite invariant family
of this shape remains separating at depth is open (OPEN\_PROBLEMS \#4).

\end{document}
